\section{Semântica funcional}

\subsection{Operações inteiras elemento a elemento}

O conjunto de operações aritméticas e lógicas considerado nesta etapa é:

\begin{itemize}
    \item soma vetorial;
    \item subtração vetorial;
    \item AND bit a bit;
    \item OR bit a bit;
    \item XOR bit a bit.
\end{itemize}

Para uma operação binária com destino |$vd$| e fontes |$vs1$| e |$vs2$|,
cada elemento ativo |$i$|, em que |$0 \leq i < vl$|, é calculado como:

|$$
vd[i] = f(vs1[i],vs2[i]).
$$|

O resultado é truncado para SEW bits. Para soma e subtração, isso
corresponde a aritmética modular:

|$$
vd[i] = f(vs1[i],vs2[i]) \bmod 2^{SEW}.
$$|

\subsection{Política de cauda}

A política definida para esta versão é tail-undisturbed. Elementos do
destino com índice |$i \geq vl$| preservam seu valor anterior.

|$$
vd_{\text{new}}[i] = vd_{\text{old}}[i], \qquad i \geq vl.
$$|

\subsection{Política de máscara}

A interface funcional aceita, quando aplicável, uma máscara representada
por uma sequência de valores booleanos. Para um elemento ativo dentro de
\texttt{vl}:

|$$
vd_{\text{new}}[i] =
\begin{cases}
f(vs1[i],vs2[i]), & mask[i] = 1, \\
vd_{\text{old}}[i], & mask[i] = 0.
\end{cases}
$$|

A política é mask-undisturbed. A validação de tamanho da máscara deve
garantir que todos os índices consultados existam. O formato arquitetural
final do registrador de máscara ainda não foi especificado.

\subsection{Loads e stores contíguos}

O modelo inclui cargas e stores unit-stride. Para um elemento de largura
SEW:

|$$
element\_bytes = \frac{SEW}{8}.
$$|

O endereço do elemento |$i$| é:

|$$
address(i) = base + i \times element\_bytes.
$$|

A memória é endereçada por byte e usa little-endian. Loads respeitam
\texttt{vl}, máscara e política de cauda. Stores escrevem apenas os
elementos ativos, isto é, aqueles com |$0 \leq i < vl$| e máscara ativa.

O nome atual das operações de memória precisa ser consolidado. O modelo
foi generalizado para operações dependentes do SEW, mas o executor e
alguns testes ainda usam nomes \texttt{vle32.v} e \texttt{vse32.v}.
Esses nomes sugerem EEW fixo de 32 bits e podem não corresponder à
largura efetivamente usada pelo modelo.